% Pursuing Stacks (cotes 134-2 to 134-8): archive cote 134-7, site page 29 (PDF page 30), typescript page 471.
% typescript. Draft compared: ps6.tex:1342-1405.
% AI-transcribed by Claude Opus 5.5 (claude-opus-5-5, adaptive thinking, effort high), Message Batches API, batch ps3 (msgbatch_01XdFmuVVBAcgc7EFxdtCqsj), 2026-11-12; not yet checked by eye.
\scanpage{134-7/29}
\mspage{471}

$W(M)(k') = M\otimes_k k'$ . If henceforth we denote by $\hand{\pi}_1(X_*)$ this group
object of $U(k)$, the relevant formula now is
\[ \text{(G)} \qquad \hand{\pi}_1(X_*(k'))\simeq \hand{\pi}_1(X_*)(k') \qquad , \]
a group isomorphism functorial with respect to variable k-algebra k',
which will imply the corresponding isomorphism
\[ \text{(G')} \qquad \hand{\pi}_1(X_*\otimes_k k')\simeq \hand{\pi}_1(X_*)\otimes_k k' \]
of group objects in $U(k')$, i.e. formation of the ``schematic'' $\hand{\pi}_1$ is compa-%
tible with ground ring extension $k \quad k'$ (provided \over{it}{$\pi_1$} exists\deleted{\unsure{x}}, for a given
$X_*$ ).

We will expect \xout{that} the map of passage to quotient
\[ \text{(H)} \qquad X_1\hand{\to}\hand{\pi}_1(X_*) = G \]
to be ``epimorphic'' in a very strong sense, stronger even than just in
the sense of presheaves, the first thought that comes to mind here is
that it should be a``\uline{submersion}'' , in the sense suggested in yesterday's
reflections in connection with the description of the Kan condition. If
however we want to be able to get for $\hand{\pi}_1(X_*(k)) = \hand{\pi}_1(X_*(k))$ \uline{any} abelian
group given beforehand, in the case $k=\underline{Z}$ say, without demanding that it
be a projective k-module, and still get it via an $X_*$ with \uline{smooth}
components, this shows that when defining a notion of ``submersion'' for
objects of $U(k)$ which may not be smooth, we should not be quite as deman-%
ding \deleted{in the definition} as suggested yesterday (cf. page 463), but find
\deleted{one} \inserted{a definition} which will include also any map $X \hand{\to} Y$ coming from an epimorphism
$M\hand{\to} N$ of k-modules (which will allow us to take $X_1$ as associated to
a projective k-module admitting the given $\hand{\pi}_1$ as it's quotient). \unsure{One}
idea that comes to mind here, is to take this property as the \uline{definition}
of a submersion, as an arrow in $U(k)$ which \xout{comes from} is isomorphic to
one obtained from an epimorphism in $Ab_k$ . This\hand{,} \deleted{at present,} of the
three definitions that have come to my mind so far for this notion, is
the one which looks the most convincing to me. I would'nt expect too
much from a complex $X_*$ , even a smooth one and satisfying the Kan condi-%
tion, unless (in terms of the three boundary maps from $X_2$ to $X_1$) it
does'nt give rise, as just explained, to a group object $G = \hand{\pi}_1(X_*)$
in $U(k)$, together with a \uline{submersion} (H). Thus definitely, when de-%
fining a schematic model category $\hand{\mathcal{M}}_n(k)$ of n-connected ss complexes
of unipotent bundles over k, I feel like insisting in case n=0 at least
upon this extra condition (plus of course $X_o = e$). The still more strin-%
gent condition one may think of, in order to have schematic models as

\margin{29}
